\DocumentMetadata{
  pdfversion=2.0,
  pdfstandard=ua-2,
  lang=en-US,
  tagging=on,
  tagging-setup={math/setup=mathml-SE,tabsorder=structure}
}
\documentclass[11pt,letterpaper]{article}
\usepackage[margin=0.68in,includefoot]{geometry}
\usepackage{newcomputermodern}
\usepackage{amsmath,amscd,mathtools}
\usepackage{tikz,adjustbox}
\providecommand{\square}{\mdlgwhtsquare}
\usepackage[hidelinks]{hyperref}
\hypersetup{
  pdftitle={Algebra First-Year Exam, January 2017, Part 1},
  pdfauthor={University of Florida Department of Mathematics},
  pdfsubject={Graduate mathematics examination},
  pdfdisplaydoctitle=true,
  bookmarksopen=true
}
\setcounter{secnumdepth}{0}
\setlength{\parindent}{0pt}
\setlength{\parskip}{0.28em}
\setlength{\emergencystretch}{3em}
\setlength{\leftmargini}{2.15em}
\setlength{\leftmarginii}{2.65em}
\setlength{\leftmarginiii}{3em}
\renewcommand{\labelenumi}{\textbf{\theenumi.}}
\pagestyle{empty}
\raggedbottom
\allowdisplaybreaks
\newenvironment{examproblems}[1]{%
  \begin{enumerate}%
  \setcounter{enumi}{#1}%
  \setlength{\topsep}{0.35em}%
  \setlength{\partopsep}{0pt}%
  \setlength{\itemsep}{0.48em}%
  \setlength{\parsep}{0.18em}%
}{\end{enumerate}}
\newenvironment{examsubparts}{%
  \begin{enumerate}%
  \setlength{\topsep}{0.18em}%
  \setlength{\partopsep}{0pt}%
  \setlength{\itemsep}{0.16em}%
  \setlength{\parsep}{0.1em}%
}{\end{enumerate}}
\newenvironment{examinstructions}{%
  \par\begingroup\itshape
}{%
  \par\endgroup\vspace{0.18em}
}
\makeatletter
\renewcommand\section{\@startsection{section}{1}{0pt}%
  {-1.25ex plus -.25ex minus -.1ex}%
  {0.55ex plus .1ex}%
  {\normalfont\large\bfseries}}
\renewcommand{\@maketitle}{%
  \newpage\null\vskip -1em%
  {\footnotesize\bfseries
    \href{https://www.ufl.edu/}{University of Florida}, \href{https://math.ufl.edu/}{Department of Mathematics}\par}%
  \vskip -0.5em%
  \tagstructbegin{tag=Title}%
    {\LARGE\bfseries\href{https://gma.math.ufl.edu/exams/algebra-first-year/}{Algebra First-Year Exam}, January 2017, Part 1\par}%
  \tagstructend%
  \vskip 0.25em%
  \tagmcbegin{artifact}%
    \hrule height 1pt%
  \tagmcend%
  \vskip 1em%
}
\makeatother
\title{Algebra First-Year Exam, January 2017, Part 1}
\author{}
\date{}
\begin{document}
\maketitle
\thispagestyle{empty}
\begin{examinstructions}
Answer four problems. Write your answers clearly in complete English sentences. You may quote results (within reason) as long as you state them clearly.
\end{examinstructions}
\begin{examproblems}{0}
\item
Let~\(G\) be a group.
\begin{examsubparts}
\item[\textnormal{(i)}]
If \(H \subseteq G\) is a normal subgroup and~\(K\) is a characteristic subgroup of~\(H\), show that~\(K\) is normal in~\(G\).
\item[\textnormal{(ii)}]
Give an example of a group and subgroups \(K \subseteq H \subseteq G\) with~\(H\) normal in~\(G\), \(K\) normal in~\(H\) but not normal in~\(G\).
\end{examsubparts}
\item
Let~\(G\) be a group.
\begin{examsubparts}
\item[\textnormal{(i)}]
Show that if~\(G\) is abelian, the subset~\(H\) of elements of finite order is a subgroup.
\item[\textnormal{(ii)}]
Give an example of a (nonabelian) group~\(G\) and elements \(x, y \in G\) such that \(xy\) has infinite order.
\end{examsubparts}
\item
Show that a group of order \(3^2 \cdot 17\) is abelian.
\item
This problem has three parts.
\begin{examsubparts}
\item[\textnormal{(i)}]
Define what it means for a group~\(G\) to be nilpotent.
\item[\textnormal{(ii)}]
Show that any nilpotent group is solvable.
\item[\textnormal{(iii)}]
Give an example of a finite solvable group that is not nilpotent.
\end{examsubparts}
\item
Let~\(G\) be a group acting on a set~\(S\), and \(H \subseteq G\) a normal subgroup. Assume that for any \(x_1, x_2 \in S\) there is a unique \(h \in H\) such that \(h(x_1)=x_2\). For \(x \in S\) let \(G_x\) be the stabilizer of~\(x\). Show that for any \(x \in S\), \(G\) is a semi-direct product of \(G_x\) and~\(H\).
\item
Show that for any \(n \geq 5\), \(A_n\) is the only nontrivial proper normal subgroup of \(S_n\).
\item
Show that the normalizer of a 5-Sylow subgroup of \(A_5\) has order 10, and is a maximal subgroup of \(A_5\).
\end{examproblems}
\end{document}
